 \documentclass[        
    a4paper,          % Tamanho da folha A4
    12pt,             % Tamanho da fonte 12pt
    chapter=TITLE,    % Todos os capitulos devem ter caixa alta
    oneside,          % Usada para impressao em apenas uma face do papel
    english,          % Hifenizacoes em ingles
    spanish,          % Hifenizacoes em espanhol
    brazilian,        % Ultimo idioma eh o idioma padrao do documento
    fleqn             % Comente esta linha se quiser centralizar as equacoes. Comente também a linha 65 abaixo
]{abntex2}

\input{lib/preambulo}

\setlength{\mathindent}{0pt} %Complementa o alinhamento de equações para totalmente a esquerda.


\trabalhoacademico{tccgraduacao}

\ehqualificacao{nao}

\removerbordasdohyperlink{sim} 

\cordohyperlink{nao}

%%%%%%%%%%%%%%%%%%%%%%%%%%%%%%%%%%%%%%%%%%%%%%%%%%%%%
%%         Informacao sobre a instituicao          %%
%%%%%%%%%%%%%%%%%%%%%%%%%%%%%%%%%%%%%%%%%%%%%%%%%%%%%

\ies{Universidade Federal do Ceará}
\iessigla{UFC}
\centro{Centro de Tecnologia}
\departamento{Departamento de Engenharia Elétrica}

%%%%%%%%%%%%%%%%%%%%%%%%%%%%%%%%%%%%%%%%%%%%%%%%%%%%%
%%        Informacao para TCC de Graduacao         %%
%%%%%%%%%%%%%%%%%%%%%%%%%%%%%%%%%%%%%%%%%%%%%%%%%%%%%

\graduacaoem{Engenharia Elétrica}
\habilitacao{bacharel} 

%%%%%%%%%%%%%%%%%%%%%%%%%%%%%%%%%%%%%%%%%%%%%%%%%%%%%
%%      Informacoes relacionadas ao trabalho       %%
%%%%%%%%%%%%%%%%%%%%%%%%%%%%%%%%%%%%%%%%%%%%%%%%%%%%%

\autor{Gabriel Gonzaga Sá Barreto}
\titulo{Título do Trabalho}
\data{2026}
\local{Fortaleza}

% Exemplo: \dataaprovacao{01 de Janeiro de 2012}
\dataaprovacao{}

%%%%%%%%%%%%%%%%%%%%%%%%%%%%%%%%%%%%%%%%%%%%%%%%%%%%%
%%           Informação sobre o Orientador         %%
%%%%%%%%%%%%%%%%%%%%%%%%%%%%%%%%%%%%%%%%%%%%%%%%%%%%%

\orientador{Prof.~Dr.~Diego de Sousa Madeira}
\orientadories{Universidade Federal do Ceará (UFC)}
\orientadorcentro{Centro de Ciências e Tecnologia (CCT)}
\orientadorfeminino{nao} % Coloque 'sim' se for do sexo feminino

%%%%%%%%%%%%%%%%%%%%%%%%%%%%%%%%%%%%%%%%%%%%%%%%%%%%%
%%          Informação sobre o Coorientador        %%
%%%%%%%%%%%%%%%%%%%%%%%%%%%%%%%%%%%%%%%%%%%%%%%%%%%%%

% Deixe o nome do coorientador em branco para remover do documento

\coorientador{}
\coorientadories{Universidade Coorientador (SIGLA)}
\coorientadorcentro{Centro do Coorientador (SIGLA)}
\coorientadorfeminino{nao} % Coloque 'sim' se for do sexo feminino

%%%%%%%%%%%%%%%%%%%%%%%%%%%%%%%%%%%%%%%%%%%%%%%%%%%%%
%%              Informação sobre a banca           %%
%%%%%%%%%%%%%%%%%%%%%%%%%%%%%%%%%%%%%%%%%%%%%%%%%%%%%

% Atenção! Deixe em branco o nome do membro da banca para remover da folha de aprovacao

% Exemplo de uso:
% \membrodabancadois{Prof. Dr. Fulano de Tal}
% \membrodabancadoisies{Universidade Federal do Ceará - UFC}


% \membrodabancadois{Prof.~Dr.~Xxxxxxx Xxxxxx Xxxxxxx}
% \membrodabancadoiscentro{Faculdade de Filosofia Dom Aureliano Matos (FAFIDAM)}
% \membrodabancadoisies{Universidade do Membro da Banca Dois (SIGLA)}
% \membrodabancatres{Prof.~Dr.~Xxxxxxx Xxxxxx Xxxxxxx}
% \membrodabancatrescentro{Centro de Ciências e Tecnologia (CCT)}
% \membrodabancatresies{Universidade do Membro da Banca Três (SIGLA)}
% \membrodabancaquatro{Prof.~Dr.~Xxxxxxx Xxxxxx Xxxxxxx}
% \membrodabancaquatrocentro{Centro de Ciências e Tecnologia (CCT)}
% \membrodabancaquatroies{Universidade do Membro da Banca Quatro (SIGLA)}
% \membrodabancacinco{Prof.~Dr.~Xxxxxxx Xxxxxx Xxxxxxx}
% \membrodabancacincocentro{Teste}
% \membrodabancacincoies{Universidade do Membro da Banca Cinco (SIGLA)}
% \membrodabancaseis{Prof.~Dr.~Xxxxxxx Xxxxxx Xxxxxxx}
% \membrodabancaseiscentro{}
% \membrodabancaseisies{Universidade do Membro da Banca Seis (SIGLA)}

\begin{document}	
	% Elementos pré-textuais
	\imprimircapa{}
	\imprimirfolhaderosto{}
	% \imprimirerrata{1-pre-textuais/errata}
	% \imprimirfolhadeaprovacao{}
	% \imprimirdedicatoria{1-pre-textuais/dedicatoria}
	% \imprimiragradecimentos{1-pre-textuais/agradecimentos}
	% \imprimirepigrafe{1-pre-textuais/epigrafe}
	% \imprimirresumo{1-pre-textuais/resumo}
	% \imprimirabstract{1-pre-textuais/abstract}
	% \renewcommand*\listfigurename{Lista de Figuras} %Se você comentar esta linha o título da lista fica: LISTA DE ILUSTRAÇÕES
	% \imprimirlistadeilustracoes
	% \imprimirlistadetabelas
% 	%\imprimirlistadequadros
	% \imprimirlistadealgoritmos
% 	%\imprimirlistadecodigosfonte
	% \imprimirlistadeabreviaturasesiglas
	% \imprimirlistadesimbolos{1-pre-textuais/lista-de-simbolos}   
	% \imprimirsumario
	
	\setcounter{table}{0}% Deixe este comando antes da primeira tabela.
	
% 	%Elementos textuais
	\textual
		\section{Modelagem}
\subsection{Sistemas de Coordenadas}
\subsection{Ângulos de Euler}


	% \input{2-textuais/1-introducao}
	% \input{2-textuais/2-fundamentacao-teorica}
	% \input{2-textuais/3-metodologia}
	% \input{2-textuais/4-resultados}
	% \input{2-textuais/5-conclusao}
	
% 	%Elementos pós-textuais	
	% \bibliography{3-pos-textuais/referencias}
% %	\imprimirglossario	
	% \imprimirapendices
		% Adicione aqui os apendices do seu trabalho
		% \input{3-pos-textuais/apendices/apendice-a}
		% \input{3-pos-textuais/apendices/apendice-b}
		% \input{3-pos-textuais/apendices/apendice-c}
		% \input{3-pos-textuais/apendices/apendice-d}
% 	\imprimiranexos
% % 		% Adicione aqui os anexos do seu trabalho
% 		\input{3-pos-textuais/anexos/anexo-a}
% 		\input{3-pos-textuais/anexos/anexo-b}		
	% \imprimirindice

\end{document}